% Generated by publication/build_sections.py; do not edit by hand.
\paragraph{Informal verdict.}
\textbf{A --- Complete solution}, with \textbf{high confidence} in both requested subparts.

\paragraph{Lean coverage (partial).}
Lean proves d\(4\)=sqrt(8/3), d\(6\)=sqrt\(2\), explicit strict improvements over the cube and dodecahedron, and the exact icosahedral lower bound. The matching d\(12\) linear-programming upper bound and the full asymptotic limit sqrt\(n\)d\(n\)=sqrt(8pi/sqrt\(3\)) remain under continuation.

\subsection{Audited informal answer}
Q804 --- Complete solution
\begin{enumerate}
\item Restatement and convention
\end{enumerate}
Let

\[
S=S^2=\{x\in\mathbb R^3:\|x\|=1\}.
\]

For every finite set \(X\subset S\) with at least two elements, define

\[
\delta(X)=\min_{\substack{x,y\in X\\x\ne y}}\|x-y\|,
\]

where \(\|\cdot\|\) is the ordinary Euclidean distance in \(\mathbb R^3\). For \(n\ge2\), let

\[
d(n)=\max_{\substack{X\subset S\\ |X|=n}}\delta(X).
\]

The problem asks:


If there is a regular polyhedron with \(n\) vertices inscribed in \(S\), is \(d(n)\) equal to its edge length?


Determine the limit of \(\sqrt n\,d(n)\) as \(n\to\infty\). Luc Pommeret


There is no apparent typo. The only genuine convention issue is the phrase ``regular polyhedron.'' I use its standard convex meaning, namely the five Platonic solids. If star polyhedra were admitted, the word ``edge'' would no longer necessarily refer to the shortest distance between vertices.
The maximum defining \(d(n)\) exists: on the compact space \((S^2)^n\), the function

\[
(x_1,\dots,x_n)\longmapsto \min_{i<j}\|x_i-x_j\|
\]

is continuous.

\begin{enumerate}
\item Final result
\end{enumerate}
Under the convex Platonic convention:

\[
\begin{array}{c|c|c|c}
n&\text{regular polyhedron}&\text{edge length on }S^2&
\text{equal to }d(n)?\\ \hline
4&\text{tetrahedron}&\sqrt{\frac83}&\text{yes}\\[2mm]
6&\text{octahedron}&\sqrt2&\text{yes}\\[2mm]
8&\text{cube}&\sqrt{\frac43}&\text{no}\\[2mm]
12&\text{icosahedron}&\sqrt{2-\frac{2}{\sqrt5}}&\text{yes}\\[2mm]
20&\text{dodecahedron}&\sqrt{2-\frac{2\sqrt5}{3}}&\text{no}.
\end{array}
\]

Thus the answer to part \(a\), as a general assertion, is no.
For part \(b\),

\[
\boxed{\displaystyle
\lim_{n\to\infty}\sqrt n\,d(n)
=
\sqrt{\frac{8\pi}{\sqrt3}}
}
\]

and numerically

\[
\sqrt{\frac{8\pi}{\sqrt3}}=3.809251227\ldots.
\]


Part \(a\): the regular polyhedra
\begin{enumerate}
\item Inner-product formulation
\end{enumerate}
For a finite \(X\subset S^2\), put

\[
s(X)=\max_{\substack{x,y\in X\\x\ne y}}\langle x,y\rangle .
\]

Since

\[
\|x-y\|^2=2-2\langle x,y\rangle
\]

for unit vectors,

\[
\delta(X)^2=2-2s(X).
\]

Thus maximizing the minimum distance is equivalent to minimizing the largest inner product.
We first record an elementary fact.
Lemma 3.1: an obtuse set has at most \(d+1\) elements
Let \(v_1,\dots,v_m\) be nonzero vectors in \(\mathbb R^d\) such that

\[
\langle v_i,v_j\rangle<0\qquad(i\ne j).
\]

Then \(m\le d+1\).
Proof
If \(m\ge d+2\), then the \(m\) vectors

\[
(v_i,1)\in\mathbb R^{d+1}
\]

are linearly dependent. Hence there are real numbers \(a_i\), not all zero, such that

\[
\sum_i a_iv_i=0,
\qquad
\sum_i a_i=0.
\]

The second equality implies that the coefficients have both signs. Set

\[
P=\{i:a_i>0\},\qquad N=\{j:a_j<0\},\qquad b_j=-a_j>0.
\]

Then

\[
w:=\sum_{i\in P}a_iv_i=\sum_{j\in N}b_jv_j.
\]

Taking the scalar product of the two expressions for \(w\),

\[
\|w\|^2
=
\sum_{\substack{i\in P\\j\in N}}
a_i b_j\langle v_i,v_j\rangle<0,
\]

a contradiction. \(\square\)

\begin{enumerate}
\item The tetrahedron: \(d(4)=\sqrt{8/3}\)
\end{enumerate}
Let \(x_1,\dots,x_4\in S^2\). Then

\[
0\le \left\|\sum_{i=1}^4x_i\right\|^2
=
4+2\sum_{1\le i<j\le4}\langle x_i,x_j\rangle.
\]

Consequently,

\[
\sum_{i<j}\langle x_i,x_j\rangle\ge-2.
\]

There are six pairs, so at least one pair satisfies

\[
\langle x_i,x_j\rangle\ge-\frac13.
\]

Thus

\[
s(X)\ge-\frac13
\]

and

\[
\delta(X)^2\le2+\frac23=\frac83.
\]

Equality is attained by

\[
\frac1{\sqrt3}(1,1,1),\quad
\frac1{\sqrt3}(1,-1,-1),\quad
\frac1{\sqrt3}(-1,1,-1),\quad
\frac1{\sqrt3}(-1,-1,1),
\]

whose distinct pairwise inner products are all \(-1/3\). These are the vertices of a regular tetrahedron.
Hence

\[
\boxed{d(4)=\sqrt{\frac83}}.
\]


\begin{enumerate}
\item The octahedron: \(d(6)=\sqrt2\)
\end{enumerate}
For six points \(x_1,\dots,x_6\in S^2\), Lemma 3.1 shows that they cannot all have negative pairwise inner products, since an obtuse set in \(\mathbb R^3\) has at most four elements. Therefore

\[
s(X)\ge0,
\]

and hence

\[
\delta(X)^2\le2.
\]

The six points

\[
\{\pm e_1,\pm e_2,\pm e_3\}
\]

have pairwise inner products \(0\) or \(-1\). Thus their largest distinct inner product is \(0\). They form a regular octahedron, and

\[
\boxed{d(6)=\sqrt2}.
\]


\begin{enumerate}
\item The cube is not optimal for \(n=8\)
\end{enumerate}
The vertices of the inscribed cube are

\[
\frac1{\sqrt3}(\pm1,\pm1,\pm1).
\]

The largest inner product between distinct vertices is \(1/3\), attained by adjacent vertices. Its edge length is therefore

\[
\sqrt{2-\frac23}=\sqrt{\frac43}.
\]

We now give a strictly better eight-point configuration.
Set

\[
u=\frac1{1+2\sqrt2}=\frac{2\sqrt2-1}{7},
\qquad
h=\sqrt u,
\qquad
r=\sqrt{1-u}.
\]

For \(k=0,1,2,3\), define

\[
A_k=
\left(
r\cos\frac{k\pi}{2},
r\sin\frac{k\pi}{2},
h
\right)
\]

and

\[
B_k=
\left(
r\cos\left(\frac\pi4+\frac{k\pi}{2}\right),
r\sin\left(\frac\pi4+\frac{k\pi}{2}\right),
-h
\right).
\]

These are the vertices of a square antiprism on \(S^2\).
Within either square, the largest distinct inner product is

\[
h^2+r^2\cos\frac\pi2=u.
\]

Between the two squares, the largest inner product occurs when the azimuthal difference is \(\pi/4\):

\[
r^2\cos\frac\pi4-h^2
=
\frac{1-u}{\sqrt2}-u.
\]

The definition \(1-u=2\sqrt2\,u\) gives

\[
\frac{1-u}{\sqrt2}-u=u.
\]

All other inner products are smaller. Hence this configuration has

\[
s(X)=u.
\]

Since

\[
u=\frac1{1+2\sqrt2}<\frac13,
\]

we obtain

\[
d(8)\ge\sqrt{2-2u}>\sqrt{2-\frac23}
=\sqrt{\frac43}.
\]

Therefore

\[
\boxed{\text{the inscribed cube is not optimal}.}
\]

Numerically, the square-antiprism construction gives minimum distance

\[
\sqrt{2-\frac{2}{1+2\sqrt2}}
=1.2155625\ldots,
\]

whereas the cube edge has length \(1.1547005\ldots\).

\begin{enumerate}
\item The icosahedron: exact optimality for \(n=12\)
\end{enumerate}
We use a self-contained linear-programming bound on the sphere.
Lemma 7.1: spherical polynomial bound
Let \(P_k\) be the degree-\(k\) Legendre polynomial, normalized by \(P_k(1)=1\). Suppose

\[
f(t)=\sum_{k=0}^m f_kP_k(t),
\]

where

\[
f_0>0,\qquad f_k\ge0\quad(k\ge1).
\]

Suppose also that

\[
f(t)\le0\qquad(-1\le t\le s).
\]

Then every \(N\)-point set \(X\subset S^2\) satisfying

\[
\langle x,y\rangle\le s\qquad(x\ne y)
\]

obeys

\[
N\le\frac{f(1)}{f_0}.
\]

Proof
The spherical-harmonic addition formula gives

\[
P_k(\langle x,y\rangle)
=
\frac{4\pi}{2k+1}
\sum_{\ell=-k}^kY_{k,\ell}(x)Y_{k,\ell}(y)
\]

for any real orthonormal basis of degree-\(k\) spherical harmonics. Consequently,

\[
\sum_{x,y\in X}P_k(\langle x,y\rangle)
=
\frac{4\pi}{2k+1}
\sum_{\ell=-k}^k
\left(\sum_{x\in X}Y_{k,\ell}(x)\right)^2
\ge0.
\]

Therefore

\[
\sum_{x,y\in X}f(\langle x,y\rangle)\ge f_0N^2.
\]

On the other hand, every off-diagonal term is nonpositive, so

\[
\sum_{x,y\in X}f(\langle x,y\rangle)\le Nf(1).
\]

Thus \(f_0N^2\le Nf(1)\), proving the result. \(\square\)
Application to twelve points
Put

\[
a=\frac1{\sqrt5},
\qquad
f(t)=(t+1)(t+a)^2(t-a).
\]

For \(-1\le t\le a\),

\[
t+1\ge0,\qquad (t+a)^2\ge0,\qquad t-a\le0,
\]

so \(f(t)\le0\).
Using

\[
\begin{aligned}
P_0(t)&=1,\\
P_1(t)&=t,\\
P_2(t)&=\frac{3t^2-1}{2},\\
P_3(t)&=\frac{5t^3-3t}{2},\\
P_4(t)&=\frac{35t^4-30t^2+3}{8},
\end{aligned}
\]

one obtains

\[
f=\sum_{k=0}^4f_kP_k
\]

with

\[
\begin{aligned}
f_0&=\frac{2(5+\sqrt5)}{75},\\
f_1&=\frac{2(5+\sqrt5)}{25},\\
f_2&=\frac{46+14\sqrt5}{105},\\
f_3&=\frac{2(5+\sqrt5)}{25},\\
f_4&=\frac8{35}.
\end{aligned}
\]

All coefficients are positive, and direct calculation gives

\[
f(1)=12f_0.
\]

The lemma therefore shows that a code with all distinct inner products at most \(1/\sqrt5\) has at most twelve points.
We must still rule out twelve points having all inner products strictly less than \(1/\sqrt5\). Suppose such a configuration existed. For \(N=12\), both inequalities in the proof of Lemma 7.1 would be equalities:

\[
f_0N^2
=
\sum_{i,j}f(\langle x_i,x_j\rangle)
=
Nf(1).
\]

Since every off-diagonal summand is nonpositive, each one must vanish. The roots of \(f\) below \(a\) are

\[
-1,\qquad -a.
\]

Thus every distinct pair would have negative inner product, contradicting Lemma 3.1.
Hence every twelve-point set satisfies

\[
s(X)\ge\frac1{\sqrt5}.
\]

Now let

\[
\varphi=\frac{1+\sqrt5}{2},
\qquad
R=\sqrt{1+\varphi^2}.
\]

The twelve normalized vectors

\[
\frac1R(0,\pm1,\pm\varphi),\qquad
\frac1R(\pm1,\pm\varphi,0),\qquad
\frac1R(\pm\varphi,0,\pm1)
\]

are the vertices of a regular icosahedron. Their distinct pairwise inner products belong to

\[
\left\{-1,-\frac1{\sqrt5},\frac1{\sqrt5}\right\}.
\]

Thus their largest inner product is \(1/\sqrt5\), and

\[
\boxed{d(12)=\sqrt{2-\frac2{\sqrt5}}}.
\]

So the regular icosahedron is optimal.

\begin{enumerate}
\item The dodecahedron is not optimal for \(n=20\)
\end{enumerate}
Let

\[
\varphi=\frac{1+\sqrt5}{2}.
\]

A standard realization of the dodecahedron is

\[
\frac1{\sqrt3}\left(
\begin{array}{c}
(\pm1,\pm1,\pm1),\\
(0,\pm\varphi^{-1},\pm\varphi),\\
(\pm\varphi^{-1},\pm\varphi,0),\\
(\pm\varphi,0,\pm\varphi^{-1})
\end{array}
\right),
\]

with the signs chosen independently in each row.
Using

\[
\varphi^2=\varphi+1,
\qquad
\varphi+\varphi^{-1}=\sqrt5,
\]

a direct calculation shows that the distinct pairwise inner products are

\[
-1,\quad-\frac{\sqrt5}{3},\quad-\frac13,\quad
\frac13,\quad\frac{\sqrt5}{3}.
\]

Thus the largest inner product is \(\sqrt5/3\), and the dodecahedral edge length is

\[
\sqrt{2-\frac{2\sqrt5}{3}}.
\]

We construct twenty points with a smaller maximum inner product.
Let

\[
N=(0,0,1),\qquad S=(0,0,-1).
\]

For \(k=0,\dots,5\), put

\[
E_k=
\left(
\cos\frac{k\pi}{3},
\sin\frac{k\pi}{3},
0
\right),
\]

and

\[
U_k=
\left(
\sqrt{\frac23}\cos\left(\frac\pi6+\frac{k\pi}{3}\right),
\sqrt{\frac23}\sin\left(\frac\pi6+\frac{k\pi}{3}\right),
\frac1{\sqrt3}
\right),
\]


\[
L_k=
\left(
\sqrt{\frac23}\cos\left(\frac\pi6+\frac{k\pi}{3}\right),
\sqrt{\frac23}\sin\left(\frac\pi6+\frac{k\pi}{3}\right),
-\frac1{\sqrt3}
\right).
\]

The set

\[
X=\{N,S\}\cup\{E_k,U_k,L_k:0\le k\le5\}
\]

has twenty elements.
The relevant maximum inner products are:

\[
\begin{array}{c|c}
\text{pair type}&\text{maximum inner product}\\ \hline
E_j,E_k&\frac12\\[1mm]
U_j,U_k\text{ or }L_j,L_k&\frac23\\[1mm]
E_j,U_k\text{ or }E_j,L_k&
\sqrt{\frac23}\cos\frac\pi6=\frac1{\sqrt2}\\[1mm]
U_j,L_k&\frac13\\[1mm]
N,U_k\text{ or }S,L_k&\frac1{\sqrt3}.
\end{array}
\]

All remaining types give no larger value. Therefore

\[
s(X)=\frac1{\sqrt2},
\qquad
\delta(X)=\sqrt{2-\sqrt2}.
\]

But

\[
\frac1{\sqrt2}<\frac{\sqrt5}{3},
\]

because \(1/2<5/9\). Hence

\[
\sqrt{2-\sqrt2}
%
\sqrt{2-\frac{2\sqrt5}{3}}.
\]

Thus

\[
\boxed{\text{the inscribed regular dodecahedron is not optimal}.}
\]


Part \(b\): the asymptotic limit
\begin{enumerate}
\item Angular separation
\end{enumerate}
Let

\[
\rho(x,y)=\arccos\langle x,y\rangle
\]

be geodesic distance on the unit sphere, and define

\[
\theta_n
=
\max_{\substack{X\subset S^2\\|X|=n}}
\min_{x\ne y}\rho(x,y).
\]

Because

\[
\|x-y\|=2\sin\frac{\rho(x,y)}2,
\]

we have

\[
d(n)=2\sin\frac{\theta_n}{2}.
\]

The caps of angular radius \(\theta_n/2\) centered at an optimal set of \(n\) points have disjoint interiors. Since such a cap has area

\[
2\pi\left(1-\cos\frac{\theta_n}{2}\right),
\]

we obtain

\[
n\,2\pi\left(1-\cos\frac{\theta_n}{2}\right)\le4\pi.
\]

In particular,

\[
\theta_n\longrightarrow0.
\]

Consequently,

\[
\frac{d(n)}{\theta_n}
=
\frac{2\sin(\theta_n/2)}{\theta_n}
\longrightarrow1.
\]

It remains to determine the asymptotics of \(\theta_n\).

\begin{enumerate}
\item The planar packing constant
\end{enumerate}
We first prove the needed planar result rather than invoking it as a black box.
Lemma 10.1: density of a periodic separated set
Let \(Y\subset\mathbb R^2\) be periodic under a lattice \(\Lambda\), with \(q\) points modulo \(\Lambda\), and suppose distinct points of \(Y\) are at distance at least \(1\). If \(F\) is a fundamental parallelogram of \(\Lambda\), then

\[
\frac{q}{|F|}\le\frac2{\sqrt3}.
\]

Proof
Consider the Voronoi tessellation of \(Y\) modulo \(\Lambda\), hence on the flat torus \(\mathbb R^2/\Lambda\). Let the \(q\) Voronoi cells have \(k_1,\dots,k_q\) sides.
Every Voronoi side is the perpendicular bisector of two sites at distance at least \(1\). Hence each cell contains the disk of radius \(1/2\) centered at its site.
A convex \(k\)-gon containing a disk of radius \(r\) centered at an interior point has area at least

\[
a_r(k)=kr^2\tan\frac{\pi}{k}.
\]

Indeed, move its supporting lines inward until tangent to the disk. If the consecutive outward normals have angular gaps \(\alpha_1,\dots,\alpha_m\), with \(m\le k\) and \(\sum\alpha_i=2\pi\), the resulting tangential polygon has area

\[
r^2\sum_{i=1}^m\tan\frac{\alpha_i}{2}
\ge
mr^2\tan\frac{\pi}{m}
\ge
kr^2\tan\frac{\pi}{k}.
\]

For \(r=1/2\), put

\[
a(k)=\frac{k}{4}\tan\frac{\pi}{k}.
\]

As a function of real \(x\ge3\),

\[
a'(x)
=
\frac14\left(
\tan\frac\pi x-\frac\pi x\sec^2\frac\pi x
\right)<0,
\]

and

\[
a''(x)
=
\frac{\pi^2}{2x^3}
\sec^2\frac\pi x
\tan\frac\pi x>0.
\]

Thus \(a\) is decreasing and convex.
Let \(V,E,F=q\) be the numbers of vertices, edges and cells in the torus tessellation. Every vertex has degree at least three, so

\[
3V\le2E.
\]

Euler's relation on the torus gives

\[
V-E+q=0.
\]

Hence \(E\le3q\), and therefore

\[
\sum_{i=1}^qk_i=2E\le6q.
\]

By convexity and monotonicity of \(a\),

\[
\begin{aligned}
|F|
&\ge \sum_{i=1}^q a(k_i)\\
&\ge q\,a\left(\frac1q\sum_i k_i\right)\\
&\ge q\,a(6)
=q\frac{\sqrt3}{2}.
\end{aligned}
\]

Thus \(q/|F|\le2/\sqrt3\). \(\square\)
Lemma 10.2: separated points in a planar domain
Let \(\Omega\subset\mathbb R^2\) be bounded, measurable, and have boundary of area zero. Let \(M_\Omega(r)\) be the largest cardinality of a subset of \(\Omega\) whose mutual distances are at least \(r\). Then

\[
\boxed{\displaystyle
\lim_{r\downarrow0}r^2M_\Omega(r)
=
\frac{2}{\sqrt3}|\Omega|.
}
\]

Proof: lower bound
Take the triangular lattice of minimum spacing \(r\). Its fundamental parallelogram has area

\[
\frac{\sqrt3}{2}r^2.
\]

Averaging over all translations of this lattice, the average number of lattice points lying in \(\Omega\) equals

\[
\frac{|\Omega|}{(\sqrt3/2)r^2}
=
\frac{2|\Omega|}{\sqrt3\,r^2}.
\]

Some translation therefore contains at least this many points in \(\Omega\). Hence

\[
\liminf_{r\downarrow0}r^2M_\Omega(r)
\ge
\frac{2}{\sqrt3}|\Omega|.
\]

Proof: upper bound
First consider a square of side \(L\) with separation \(1\). Repeat the configuration periodically with period \(L+1\) in both coordinate directions. The resulting periodic set remains \(1\)-separated, and therefore Lemma 10.1 gives

\[
M_{[0,L]^2}(1)\le\frac2{\sqrt3}(L+1)^2.
\]

Now cover \(\Omega\), for arbitrary \(\varepsilon>0\), by finitely many grid squares with disjoint interiors and total area at most \(|\Omega|+\varepsilon\). Assign boundary points to one square. Applying the preceding square estimate after scaling by \(1/r\) gives

\[
\limsup_{r\downarrow0}r^2M_\Omega(r)
\le
\frac2{\sqrt3}(|\Omega|+\varepsilon).
\]

Letting \(\varepsilon\downarrow0\) proves the upper bound. \(\square\)
The coefficient \(2/\sqrt3\) is the point density of the triangular lattice with unit minimum spacing.

\begin{enumerate}
\item Transfer to a smooth surface
\end{enumerate}
We prove a slightly more general statement.
Proposition 11.1
Let \(M\) be a compact smooth Riemannian surface of area \(A\). Let \(P_M(r)\) be the maximum number of points of \(M\) with pairwise geodesic distance at least \(r\). Then

\[
\boxed{\displaystyle
\lim_{r\downarrow0}r^2P_M(r)
=
\frac{2A}{\sqrt3}.
}
\]

Proof
Fix \(\eta>0\). By using sufficiently small normal-coordinate neighborhoods and a finite geodesic subdivision of \(M\), we may partition \(M\), up to finitely many piecewise smooth boundaries, into cells

\[
Q_j=\phi_j(\Omega_j),\qquad 1\le j\le m,
\]

such that for \(u,v\in\Omega_j\),

\[
(1-\eta)|u-v|
\le
d_M(\phi_j(u),\phi_j(v))
\le
(1+\eta)|u-v|,
\]

and for every measurable \(E\subset\Omega_j\),

\[
(1-\eta)|E|
\le
\operatorname{area}(\phi_j(E))
\le
(1+\eta)|E|.
\]

These estimates follow from the continuity of the metric tensor in normal coordinates; the cells may be chosen inside strongly convex coordinate balls.
Upper bound
Let \(X\subset M\) be \(r\)-separated, and assign each boundary point to one cell. The preimages of the points of \(X\cap Q_j\) are Euclidean \(r/(1+\eta)\)-separated. Lemma 10.2 gives

\[
\limsup_{r\downarrow0}
r^2|X\cap Q_j|
\le
\frac2{\sqrt3}(1+\eta)^2|\Omega_j|.
\]

Summing over \(j\), and using

\[
\sum_j|\Omega_j|\le\frac{A}{1-\eta},
\]

we obtain

\[
\limsup_{r\downarrow0}r^2P_M(r)
\le
\frac{2A}{\sqrt3}
\frac{(1+\eta)^2}{1-\eta}.
\]

Lower bound
Remove from every \(Q_j\) a geodesic collar of width \(r/2\) along its boundary, and denote the remaining set by \(Q_{j,r}\). Its preimage \(\Omega_{j,r}\) satisfies

\[
|\Omega_{j,r}|\longrightarrow|\Omega_j|.
\]

Inside \(\Omega_{j,r}\), use a triangular lattice of spacing

\[
\frac{r}{1-\eta}.
\]

By the averaging argument from Lemma 10.2, we can choose at least

\[
\frac2{\sqrt3}
\frac{(1-\eta)^2}{r^2}
|\Omega_{j,r}|
\]

points. Their images are \(r\)-separated within \(Q_j\).
Points chosen in distinct cells are also \(r\)-separated: every path joining them must travel at least \(r/2\) from the first point to the boundary of its cell and at least \(r/2\) from the boundary of the second cell to the second point.
Thus

\[
\liminf_{r\downarrow0}r^2P_M(r)
\ge
\frac2{\sqrt3}(1-\eta)^2
\sum_j|\Omega_j|
\ge
\frac{2A}{\sqrt3}
\frac{(1-\eta)^2}{1+\eta}.
\]

Finally let \(\eta\downarrow0\). \(\square\)
For the unit sphere, \(A=4\pi\), so

\[
P_{S^2}(r)
\sim
\frac{8\pi}{\sqrt3\,r^2}.
\]


\begin{enumerate}
\item From the packing count to \(\theta_n\)
\end{enumerate}
Set

\[
C=\frac{8\pi}{\sqrt3}.
\]

We have

\[
r^2P_{S^2}(r)\longrightarrow C.
\]

By definition,

\[
P_{S^2}(r)\ge n
\quad\Longleftrightarrow\quad
\theta_n\ge r.
\]

This asymptotic inverse relation yields

\[
n\theta_n^2\longrightarrow C.
\]

For completeness, fix \(\varepsilon>0\). For all sufficiently small \(r\),

\[
\frac{C-\varepsilon}{r^2}
\le P_{S^2}(r)
\le
\frac{C+\varepsilon}{r^2}.
\]

Since \(P_{S^2}(\theta_n)\ge n\),

\[
n\theta_n^2\le C+\varepsilon
\]

for large \(n\).
Conversely, set

\[
r_n=\sqrt{\frac{C-2\varepsilon}{n}}.
\]

For large \(n\),

\[
P_{S^2}(r_n)
\ge
\frac{C-\varepsilon}{r_n^2}
=
n\,\frac{C-\varepsilon}{C-2\varepsilon}
n.
\]

Hence \(\theta_n\ge r_n\), so

\[
n\theta_n^2\ge C-2\varepsilon.
\]

Letting \(\varepsilon\downarrow0\),

\[
\sqrt n\,\theta_n\longrightarrow\sqrt{\frac{8\pi}{\sqrt3}}.
\]

Finally, since \(d(n)/\theta_n\to1\),

\[
\boxed{\displaystyle
\lim_{n\to\infty}\sqrt n\,d(n)
=
\sqrt{\frac{8\pi}{\sqrt3}}.
}
\]

This agrees with the general best-packing formula for rectifiable sets, whose specialization to \(S^2\) gives the same constant. arXiv

\begin{enumerate}
\item Boundary checks and normalization
\end{enumerate}
For comparison:

\[
d(2)=2
\]

using antipodal points, and

\[
d(3)=\sqrt3
\]

using an equilateral triangle on a great circle. Indeed,

\[
0\le\|x_1+x_2+x_3\|^2
\]

implies that some pair has inner product at least \(-1/2\).
The radius-one hypothesis fixes the scale. On a sphere of radius \(R\), all chord distances are multiplied by \(R\), and the limiting constant becomes

\[
R\sqrt{\frac{8\pi}{\sqrt3}}.
\]

The distinction between chord and angular distance is essential for the exact finite values, but not for the leading asymptotic term, because

\[
2\sin(\theta/2)\sim\theta
\qquad(\theta\to0).
\]
